\section*{Chapitre \ref{ch:b2:particle-diffusion} --- Diffusion de particules}

\begin{solution}{exo:b2:particle-diffusion:1}
$L^2/D$~: sucre $5 \times 10^6\,\unit{s}$ (deux mois)~; parfum de même~;
tissu \qty{500}{s}~; cellule \qty{0.05}{s}. La diffusion fait réellement le dernier
(et, tout juste, l'avant-dernier)~; la convection fait le reste.
\end{solution}

\begin{solution}{exo:b2:particle-diffusion:2}
$j = D\Delta n/e = \qty{1e-5}{mol/m^2/s}$~; $\times \qty{1e-4}{m^2}$~:
\qty{1e-9}{mol/s}, $6 \times 10^{14}$ molécules par seconde~;
$R_{\text{d}} = e/DS = \qty{1e9}{s/m^3}$~; deux en série~: moitié moins.
\end{solution}

\begin{solution}{exo:b2:particle-diffusion:3}
$\sqrt{2Dt}$~: \qty{45}{\micro m}, \qty{2.7}{mm}, \qty{13}{mm}~; sommet $N/\sqrt{4\pi
Dt}$~: \qty{9e23}{m^{-3}}, \qty{1.5e22}{m^{-3}}, \qty{3e21}{m^{-3}}~; $1\%$
après $10^4$ fois plus longtemps~: \qty{1e4}{s}.
\end{solution}

\begin{solution}{exo:b2:particle-diffusion:4}
$D \approx \qty{1.2e-5}{m^2/s}$~; $N = (L/\ell)^2 = 2 \times 10^{14}$~; $N\ell/v^*
= \qty{3e4}{s}$ contre \qty{2}{ms}. Eau~: l'estimation donne
\qty{2e-8}{m^2/s}, dix fois les valeurs mesurées --- dans un liquide une
molécule s'agite dans la cage de ses voisines et ses pas successifs
sont anticorrélés~; le pas effectif est plus court.
\end{solution}

\begin{solution}{exo:b2:particle-diffusion:5}
(a) $Dn'' = q$~: $n = n_0 - q(a^2 - x^2)/2D$. (b) $n(0) \ge 0$~: $a \le
\sqrt{2Dn_0/q}$. (c) \qty{0.28}{mm}. (d) Aucune cellule ne peut être à plus de
quelques centaines de micromètres d'un capillaire (en pratique $50$--\qty{100}{\micro m}
dans un tissu actif).
\end{solution}

\begin{solution}{exo:b2:particle-diffusion:6}
(a) $n = n_{\text{s}}a/r$, $\Phi = 4\pi D a n_{\text{s}} = \qty{3.6e-8}{mol/s}$.
(b) $1.95 \times 10^{-5}\,\unit{mol}$~; \qty{540}{s}. (c) La couche en $1/r$ met
$\sim a^2/D \approx \qty{2000}{s}$ à s'établir et le grain rétrécit~; remuer
remplace la couche d'épaisseur $\sim a$ par une \omterm{prop:b2:viscous-flows:bl}{couche limite} $\delta \ll a$
et multiplie le flux par $a/\delta$. (d) $(rn')' = 0$~: $n = A + B\ln r$,
qui ne peut s'annuler à l'infini~: pas d'état stationnaire --- le nuage d'un
cylindre ne cesse de croître (logarithmiquement).
\end{solution}

\begin{solution}{exo:b2:particle-diffusion:7}
(a) $\dd/\dd t\int n = D[n']_{-\infty}^{\infty} = 0$~; $\dd/\dd t\int x^2 n = D\int
x^2 n'' = 2D\int n = 2DN$. (b) $\partial_t[n(x,-t)] = -D\,\partial_x^2 n$~: le
signe s'inverse. (c) Pour $t < 0$ la «~largeur~» $2Dt$ est négative~: pas de tel
état --- la bouffée ne peut être déconcentrée. (d) $\dd/\dd t\int n^2 = 2D\int nn''
= -2D\int n'^2 \le 0$.
\end{solution}

\begin{solution}{exo:b2:particle-diffusion:8}
(a) \cref{prop:b2:particle-diffusion:erfc}. (b) $x = 2 \times 1.16\sqrt{Dt}
= \qty{0.62}{mm}$. (c) $4\times$~: \qty{16}{h}. (d) $4/2.5 = \qty{1.6}{h}$.
\end{solution}

\begin{solution}{exo:b2:particle-diffusion:9}
(a) $D = \qty{2.2e-13}{m^2/s}$~; $\sqrt{2Dt}$~: \qty{0.66}{\micro m},
\qty{5}{\micro m}. (b) \qty{7e-11}{m^2/s}. (c) $a = k_BT/6\pi\eta D =
\qty{0.44}{nm}$. (d) $\langle x^2\rangle = 2Dt = RTt/3\pi\eta a N_A$~: toutes les
grandeurs sauf $N_A$ sont mesurées.
\end{solution}

\begin{solution}{exo:b2:particle-diffusion:10}
(a) $Dn_1S/e$. (b) Les particules mettent $\sim e^2/D$ à traverser. (c)
$e^2/6D = \qty{670}{s}$, onze minutes. (d) Le retard donne $D$ seul~;
le flux stationnaire donne $D \times$ (solubilité)~: deux mesures,
deux inconnues.
\end{solution}

\begin{solution}{exo:b2:particle-diffusion:11}
(a) $n = n_\infty(1 - a/r)$, $\Phi = 4\pi D a n_\infty$. (b) $3.8 \times
10^6$ par seconde. (c) Les récepteurs agissent comme des conductances en parallèle puis
en série avec la couche sphérique~: $\Phi \approx 4\pi Dan_\infty
\cdot N_{\text{r}}s/(N_{\text{r}}s + \pi a)$~; la moitié du maximum pour
$N_{\text{r}}s = \pi a$~: avec $s = \qty{1}{nm}$, $N_{\text{r}} \approx
3000$, couvrant $N_{\text{r}}s^2/4a^2 \approx 10^{-3}$ de la surface.
(d) Chaque sorte de récepteur ne demande qu'une aire négligeable pour une capture
quasi maximale~: une cellule peut surveiller des milliers de substances à la fois.
\end{solution}

\begin{solution}{exo:b2:particle-diffusion:12}
(a) Différence avant en $t$, différence seconde centrée en $x$. (b)
$\alpha = 1/2$~: $n_i^{k+1} = (n_{i-1}^k + n_{i+1}^k)/2$ --- chaque particule saute
à gauche ou à droite avec la probabilité $1/2$. (c) Pour $n_i = (-1)^i$, $n^{k+1} =
(1 - 4\alpha)n^k$~: $|1 - 4\alpha| > 1$ quand $\alpha > 1/2$. (d) $\Delta t \le
\Delta x^2/2D = \qty{3.3}{s}$~; environ $1100$ pas.
\end{solution}

\begin{solution}{pb:b2:particle-diffusion:1}
\textbf{1.} $k_BT = \qty{0.118}{eV}$, $E_{\text{a}}/k_BT = 29.2$~: $D =
\qty{1.6e-17}{m^2/s}$~; à \qty{1000}{\celsius}~: \qty{1.6e-18}{m^2/s} --- un
facteur $10$.

\textbf{2.} La surface réfléchissante se traite par l'image miroir de la
bouffée, ce qui double l'amplitude sur $x > 0$~; $\int_0^\infty = (Q/\sqrt{\pi Dt})
\cdot \tfrac12\sqrt{4\pi Dt} = Q$.

\textbf{3.} $Dt = \qty{5.8e-14}{m^2}$~; $\sqrt{\pi Dt} = \qty{0.43}{\micro m}$~; $n_{\text{s}}
= \qty{2.3e24}{m^{-3}}$.

\textbf{4.} $\eu^{-x^2/4Dt} = 4.3 \times 10^{-4}$~: $x_j^2 = 4Dt \times 7.75$, $x_j =
\qty{1.3}{\micro m}$.

\textbf{5.} $n_{\text{s}}$ est divisé par deux, le logarithme tombe à $7.05$, $4Dt$
quadruple~: $x_j = \qty{2.6}{\micro m}$ --- moins que doublé.

\textbf{6.} $\delta D/D = (E_{\text{a}}/k_BT)\,\delta T/T = 21\%$~; $x_j \propto
\sqrt{Dt}$ en gros, donc $1\%$ sur $x_j$ demande $2\%$ sur $D$~: \qty{1}{K}.

\textbf{7.} $E_{\text{a}}/k_BT = 134$~: $D \approx 10^{-62}\,\unit{m^2/s}$~; une
distance interatomique (\qty{0.25}{nm}) en $a^2/D \sim 10^{43}\,\unit{s}$~: gelé.

\textbf{8.} $E_{\text{a}}/k_BT = 32.8$~: $D = \qty{4.3e-19}{m^2/s}$~; $\sqrt{Dt} =
\qty{28}{nm}$.

\textbf{9.} $n_0\operatorname{erfc}(x/2\sqrt{Dt})$~; $n_B/n_0 = 5 \times 10^{-6}$~:
$x \approx 6.4\sqrt{Dt} = \qty{0.18}{\micro m}$.

\textbf{10.} $2n_0\sqrt{Dt/\pi} = \qty{6.3e18}{m^{-2}}$~: six fois $Q$ ---
la dose croît en $\sqrt t$ et est fixée par la durée.

\textbf{11.} Basse température~: $D$ petit, dose contrôlée par la durée,
peu profonde~; haute température~: la dose fixée est enfoncée vite et loin.

\textbf{12.} \qty{0.18}{\micro m} contre \qty{1.3}{\micro m}~: mince.

\textbf{13.} $4 \times 10^{18}$ molécules~; $n_{\text{a}} = \qty{2.4e25}{m^{-3}}$.

\textbf{14.} $\ell = \qty{37}{nm}$~; $v^* = \qty{206}{m/s}$.

\textbf{15.} $D \approx \qty{2.5e-6}{m^2/s}$.

\textbf{16.} $L^2/D = 10^7\,\unit{s}$, des mois~; avec dérive~: \qty{50}{s}.

\textbf{17.} $n = 8N(4\pi Dt)^{-3/2}\eu^{-r^2/4Dt}$~; maximal en $t^* = r^2/6D =
\qty{6.7e4}{s}$, dix-huit heures.

\textbf{18.} $n_{\max} = 8N(4\pi Dt^*)^{-3/2}\eu^{-3/2} \approx \qty{2e18}{m^{-3}}$,
bien au-dessus du seuil~: senti, à la longue.

\textbf{19.} $r^2 = 4Dt\ln 10^6 = \qty{12}{m^2}$~: \qty{3.5}{m}.

\textbf{20.} $v^*t/\ell \approx 5 \times 10^{14}$.

\textbf{21.} $\langle x\rangle = 0$, $\langle x^2\rangle = N\ell^2 = \ell^2 t/\tau$~;
$D = \ell^2/2\tau$.

\textbf{22.} $\binom{N}{N/2}/2^N$~: $1/2$, $3/8$, $5/16$ --- décroissante
(comme $1/\sqrt{\pi N/2}$)~: le marcheur s'éloigne en $\sqrt N$.

\textbf{23.} Le nuage se rassemble en un point --- jamais vu~; cela
viole l'\omterm{thm:b2:particle-diffusion:equation}{équation de diffusion}, non les lois de la mécanique.

\textbf{24.} Chaque collision est réversible, mais l'état initial (toutes les
particules ensemble) est exceptionnel~: presque toute histoire microscopique
issue de lui s'étale, et l'inverse exige une conspiration de toutes les
vitesses. La moyenne de la marche au hasard ne retient que ce qui est typique et
jette cette information~; l'irréversibilité est statistique.

\textbf{25.} $L \sim \sqrt{Dt}$, $t \sim L^2/D$~; $D \sim \ell v^*/3$, c'est-à-dire
le pas entre collisions~: $10^{-5}$, $10^{-9}$, $10^{-17}\,\unit{m^2/s}$.
\end{solution}
